\cotefile{31}
\scanpage{52}
% archive sheet number: 51

% L1–4, upper lemma
\uline{Lemme 2} Soit $k$ un corps, tel que $H^i(k',\Gamma)=0$ pour $\unsure{i|l}>n$. \inserted{et tout groupe abélien fini $\Gamma$ où $\unsure{G'_k|G_{k'}}$ opère trivialement, et tte extension finie $k'$ de $k$.} Soit $K$ une extension de \deleted{\unsure{$k$}} type fini de $k$, de degré de transcendance $1$. Alors $H^i(K,\Gamma)=0$ pour $i>n+1$.

% L5–7, left field diagram
\diagram{\begin{tikzcd}[ampersand replacement=\&,column sep=small,row sep=small]
k \arrow[d,dash] \& \&\\
k'=\widetilde{k}\cap K \arrow[d,dash] \& K \arrow[l,dash] \arrow[d,"G''"'] \arrow[dr,"G"] \&\\
\widetilde{k} \arrow[r,dash] \& \widetilde{k}K \arrow[r,"G'"'] \& \widetilde{K}\deleted{K}
\end{tikzcd}}

% L5–7, right of diagram
\unsure{Cons.|On a|d'où} \unsure{sa|la|une} suite exacte

% L6–7, exact sequence
\[
e \mathbin{\unsure{-|\to}} G' - G - G'' - e,
\]

% L8–12, spectral sequence and vanishing
d'où une suite spectrale $\underline{H}^i(G'',\underline{H}^j(G',\Gamma))\Longrightarrow\underline{H}^{\unsure{*|\times}}(G,\Gamma)$. Or $H^j(G',\Gamma)=0$ si $j>1$ en vertu du lemme 1, et d'autre part, \deleted{$G''$ étant un sous-groupe d'indice fini de} \inserted{$\underline{H}^i(G'',\underline{H}^j(G',\Gamma))=0$ si $i>n$, par} l'hypothèse. D'où par suite $H^p(G,\Gamma)=0$ si $p>n+1$, cqfd.

% L13–16, corollary
\uline{Corollaire} Soit $K$ un corps de fonctions algébriques sur un corps algt. clos $k$, $\mathrm{deg.tr.}_k K=n$. Alors $H^i(K,\Gamma)=0$ pour $i>n$ et tout groupe \inserted{abélien} fini $\Gamma$ où $G_K$ opère trivialement.
